\section{Proofs from corrected priorities to rationality}
\label{sec:annex-p18-structured-priority-rationality}

The strict comparisons are those of the corrected general account
\cite{ModgilPrakken2013GeneralCorrected}. The transposition construction
below preserves the scope of Caminada--Amgoud's assumption-retaining
Lemma 6 \cite{CaminadaAmgoud2007Evaluation}.

\begin{annexitem}{annex:p18.priorityorder}{thm:p18-priority-order}{Strict argument orders}
\begin{proof}
Each typed strict set comparison is a strict partial order by
Theorem~\ref{thm:p18-rationality-strict-sets}.
If \(A\) is strict and firm, both its ordinary-premise and defeasible-rule
sets are empty, as is its last-rule set. An empty left operand never
compares below anything. Thus \(A\) is below no argument under either
link principle.
If \(B\) is fallible, last link either compares its nonempty last-rule
set with the empty set, or, when \(B\) is strict, compares its nonempty
ordinary-premise set with the empty set. Either comparison makes
\(B\prec A\).
For weakest link, a strict plausible \(B\) is below \(A\) by its
premise set, a firm defeasible \(B\) by its rule set, and a plausible
defeasible \(B\) by both nonempty sets. Thus firm strict arguments
have the stated position.

Irreflexivity of last link follows because either strict set comparison
of a component with itself is impossible.
For transitivity suppose \(C\prec_{\rm L}B\prec_{\rm L}A\).
If \(B\) is defeasible, \(C\) is defeasible and
\(\operatorname{Last}(C)\triangleleft\operatorname{Last}(B)
\triangleleft\operatorname{Last}(A)\); transitivity proves \(C\prec A\).
If \(B\) is strict, \(A\) is strict, since an empty last-rule set
cannot compare below a nonempty one. Their ordinary-premise sets compare.
If \(C\) is defeasible it is below the strict \(A\) by the empty-right
clause. If \(C\) is strict, transitivity of the premise comparison
gives \(C\prec A\).

For weakest link, write \(P_X=\operatorname{Prem}_p(X)\) and
\(D_X=\operatorname{Def}(X)\).
Irreflexivity follows in each of the three ordered cases from
irreflexivity of the participating set comparison.
Suppose \(C\prec_{\rm W}B\prec_{\rm W}A\).
The middle argument \(B\) cannot be strict and firm.
If \(B\) is strict and plausible, then \(A\) must be strict:
otherwise the comparison would require
\(\varnothing\triangleleft D_A\).
Thus \(P_B\triangleleft P_A\).
If \(C\) is strict, transitivity gives
\(P_C\triangleleft P_A\), proving the result.
If \(C\) is defeasible, its comparison with \(B\) requires
\(P_C\triangleleft P_B\) and \(D_C\triangleleft\varnothing\);
the first gives \(P_C\triangleleft P_A\), and \(D_A=\varnothing\)
gives the second comparison for \(C\prec A\).
If \(B\) is firm and defeasible, the symmetric argument exchanges
premise and defeasible-rule sets.
Finally, if \(B\) is plausible and defeasible, both \(P_B,D_B\)
are nonempty. Then \(C\prec B\) requires both
\(P_C\triangleleft P_B\) and \(D_C\triangleleft D_B\).
Likewise \(B\prec A\) requires
\(P_B\triangleleft P_A\) and \(D_B\triangleleft D_A\), whatever
empty components \(A\) may have. Transitivity gives both comparisons
for \(C\prec A\). This proves transitivity in every case.
\end{proof}
\end{annexitem}

\begin{annexitem}{annex:p18.priorityresources}{thm:p18-priority-resources}{Minimal and maximal finite witnesses}
\begin{proof}
Every nonempty finite strict partial order has a minimal and a maximal
element. Otherwise successively choosing a strictly smaller, or larger,
element would force repetition in a finite set and then a strict
self-comparison by transitivity.

Suppose \(\Gamma\triangleleft_{\rm E}\Delta\) with \(\Delta\ne\varnothing\).
Choose \(g\in\Gamma\) below every element of \(\Delta\).
Descend from \(g\) in \(\Gamma\cup\Delta\) to a minimal element \(m\).
No element of \(\Delta\) can be \(g\) or lie strictly below \(g\):
if \(d\in\Delta\) did so, \(g<d\) and \(d\leq g\), where this last
notation means equality or strict inequality, would contradict
irreflexivity. Therefore \(m\in\Gamma\).
Minimality says that no element of \(\Gamma\cup\Delta\) is below \(m\).
If \(m\in\Sigma\) and \(H\subseteq\Gamma\cup\Delta\), no element
of \(H\) can be below every element of \(\Sigma\), because it cannot
be below \(m\). The empty-left clause covers \(H=\varnothing\).
The same argument starts at any minimal element of \(\Gamma\) when
\(\Delta=\varnothing\).

Suppose \(\Gamma\triangleleft_{\rm D}\Delta\ne\varnothing\).
Each element of \(\Gamma\) is below an element of \(\Delta\).
A maximal element \(d\) of \(\Gamma\cup\Delta\) therefore belongs
to \(\Delta\), and is below no element of \(\Gamma\).
If \(d\in H\) and \(\varnothing\ne\Sigma\subseteq\Gamma\), this
particular member of \(H\) has no strictly larger witness in \(\Sigma\);
the democratic comparison fails.
Finally if \(\varnothing\ne\Sigma\subseteq H\), take a maximal
element \(s\) of \(\Sigma\). This \(s\in H\) has no strictly larger
element in \(\Sigma\), so again \(H\not\triangleleft_{\rm D}\Sigma\).
\end{proof}
\end{annexitem}

\begin{annexitem}{annex:p18.prioritycontinuations}{thm:p18-priority-continuations}{Repeated families and finite descent}
\begin{proof}
Adding only strict nodes and axiom leaves creates no new ordinary premise
or defeasible rule. Strict nodes take unions of their children's last-rule
sets. Since each indexed input is retained, all three resulting component
sets are exactly their input unions, not merely subsets.
For a singleton input, its relevant components and its strict/firm
classification therefore remain unchanged. Both link comparisons with
every other argument are unchanged. This proves clause (iii).
Clauses (i) and (ii) are the strict-firm assertions of the order theorem.

For clause (iv), suppose, to the contrary, that every omitted-index
continuation \(C_{-i}\) is strictly below its corresponding \(C_i\).
For last link, if some \(C_i\) is defeasible, let \(I\) be the finite
nonempty set of indices with nonempty last-rule sets \(L_i\).
For each \(i\in I\), the assumed comparison necessarily gives
\[
 \bigcup_{j\ne i}L_j\triangleleft L_i.
 \tag{9}
\]
By the proved finite-union property, some \(j\ne i\) satisfies
\(L_j\triangleleft L_i\); its left operand is nonempty, so \(j\in I\).
Starting at any index in \(I\), repeatedly choose such an index \(j\).
The finite index set forces a repeated index. The resulting cycle of
strict set comparisons contradicts transitivity and irreflexivity.
This reasoning does not require distinct \(C_i\)'s or distinct \(L_i\)'s.
Indeed, equal component sets cannot themselves form a strict comparison.

If all \(C_i\) are strict, so are their strict continuations.
The assumed comparisons use the ordinary-premise sets instead.
If all those sets are empty, no strict comparison is possible.
Otherwise repeat the preceding finite-index descent on the nonempty
ordinary-premise sets. This proves clause (iv) for last link.

For weakest link, if some \(C_i\) has nonempty defeasible-rule set
\(G_i\), its comparison with \(C_{-i}\) must use the rule comparison,
either alone or in the conjunction case. Thus
\(\bigcup_{j\ne i}G_j\triangleleft G_i\).
The same finite-union witness and index-cycle argument is a contradiction.
If every \(G_i\) is empty, all the arguments and continuations are strict,
and the ordinary-premise argument applies again.
Thus weakest link also satisfies clause (iv), with repetitions allowed.

For profile continuations, ordinary-premise and defeasible-rule unions
are exactly the required equalities in their definition.
These are all the components used by weakest link, so its singleton
invariance and the same finite-index argument remain valid.
The proof does not assert equality of last-rule sets for arbitrary
profile continuations.
\end{proof}
\end{annexitem}

\begin{annexitem}{annex:p18.priorityprofile}{thm:p18-priority-profile-boundary}{A repeated rule changes the last link}
\begin{proof}
Use paired literals over \(q,r,\delta_1,\delta_2\), let \(K_n=\{q\}\), \(K_p=\varnothing\),
and let the strict-rule set be empty. Add
\[
 d_1:q\Rightarrow r,\qquad d_2:r\Rightarrow q,
 \qquad d_1<_d d_2.
 \tag{10}
\]
Set \(n(d_i)=\delta_i\); no argument concludes their negatives.
Let
\[
 Q=[q],\quad U=[Q;d_1],\quad A=[U;d_2],\quad A'=[A;d_1].
 \]
The arguments \(A,A'\) have the same ordinary-premise set
\(\varnothing\), the same axiom-premise set \(\{q\}\), the same
defeasible-rule set \(\{d_1,d_2\}\), and empty strict-rule sets.
Thus \(A'\) is a profile continuation of \(\{A\}\).
But their last-rule sets are respectively \(\{d_2\}\) and \(\{d_1\}\).
For elitism and democracy alike, \(A'\prec_{\rm L}A\).
Irreflexivity gives \(A\not\prec_{\rm L}A\), so clause (iii), applied
with its comparison argument equal to \(A\), would require
\(A'\not\prec_{\rm L}A\), a contradiction.
The new top application of \(d_1\) is defeasible, so \(A'\) is not
an actual strict continuation of \(A\).
The component-set condition loses this new application because \(d_1\)
was already present earlier. The countermodel concerns the stated
continuation condition; it does not refute the directly proved return
or rationality results for last link.
\end{proof}
\end{annexitem}

\begin{annexitem}{annex:p18.priorityreversal}{thm:p18-priority-retained-reversal}{Occurrence-sensitive transposition}
\begin{proof}
Treat the assumptions as added nullary strict rules, exactly as in the
source lemma. Fix an occurrence of the selected assumption \(a_i\).
If the whole tree is that assumption node, the added assumption
\(\neg c=\neg a_i\) gives the desired derivation.
Otherwise follow the unique path from the root to the selected occurrence.
At the root, whose rule is
\((\phi_1,\ldots,\phi_j,\ldots,\phi_m)\to c\), use its transposition
\[
 (\phi_1,\ldots,\neg c,\ldots,\phi_m)\to\neg\phi_j .
\]
Supply its other antecedents with the original sibling derivations.
This constructs a derivation of the negation of the selected child's
conclusion from the original assumptions and \(\neg c\).
Repeat at the next node on the path. Each step moves strictly down
the finite original tree, so eventually it concludes \(\neg a_i\).
All added rule applications are transpositions of original strict rules.
The original assumption \(a_i\) may still occur in a sibling derivation.
The proof has never claimed that it can be removed from the available
assumption set. This is also an induction on the path length, and
is the source's depth induction with its chosen occurrence made explicit.

For the grafted version, begin with the actual argument \(X\) as the
opposing root argument, and use the actual off-path subtrees of \(B\)
as the other children at each transposed rule. Stop on reaching
the selected \(U\in M(B)\). Its root is ordinary or defeasible, so
the final opposing conclusion undermines or rebuts at that root.
The new graft contains \(X\). Its only new nodes are strict.
Every attack on it lands at a fallible subtree already in \(X\)
or in one of the inserted subtrees of \(B\). The same target and
preference comparison then gives a defeat of \(X\) or \(B\).
This proves the stated defeater condition.
Its premise and defeasible-rule sets are contained in the corresponding
unions. Because the reversed path lies entirely above the frontier,
all off-path last-rule sets are subsets of \(\operatorname{Last}(B)\);
strict grafting takes unions of last-rule sets. This proves the last-rule
claim as well. Distinct occurrences of an identical subtree are handled
by distinct positions on the original tree, not by assuming distinct trees.
\end{proof}
\end{annexitem}

\begin{annexitem}{annex:p18.prioritycontra}{thm:p18-priority-contraposition-graft}{Two-root contraposition}
\begin{proof}
Put \(b=\operatorname{Conc}(B)\), \(x=\operatorname{Conc}(X)\),
and \(u=\operatorname{Conc}(U)\).
The premise set \(\{b,u\}\) strictly derives \(b\) by its premise leaf.
Apply closure under contraposition with the removed premise \(u\)
and the chosen contradictory \(x\) of \(b\).
It gives a finite strict derivation of some contradictory \(u^\ast\)
of \(u\) from a subset of \(\{b,x\}\).
If \(b=u\), the resulting premise set may simply be \(\{x\}\);
the same conclusion and construction apply.

Substitute \(B\) for any leaf labelled \(b\), and \(X\) for any leaf
labelled \(x\). When these labels coincide, either whole argument may
be chosen consistently at each occurrence; choose \(B\) if a single
choice is needed. Premises unused by the strict derivation require
no substitution. The result \(W\) concludes \(u^\ast\) and therefore
attacks \(B\) at its fallible frontier node \(U\).
All its fallible targets lie in the substituted whole arguments.
Consequently every defeater of \(W\) defeats \(X\) or \(B\).
If no premise is used, \(W\) has no ordinary premise or defeasible rule,
and is strict and firm.
Its premise set is a subset of the union of the two original premise
sets; joint c-consistency therefore makes \(W\) an admitted c-restricted
argument. This proof uses full contraposition, which entails the required
strict derivation from opposing conclusions; it does not infer this
derivation from rule transposition alone.
\end{proof}
\end{annexitem}

\begin{annexitem}{annex:p18.priorityreturn}{thm:p18-priority-return}{Priority witnesses produce a returning defeat}
\begin{proof}
Since \(B T A\) fails as a defeat, its target \(X\) is an ordinary
premise or a defeasible-root subtree, the conflict is contradictory,
and \(B\prec X\). Undercuts and one-way contrary attacks cannot fail.
By the strict-order theorem, \(X\not\prec B\).
The argument \(B\) cannot be strict and firm.
If its own root is ordinary or defeasible, take \(W=X\):
the contradictory attack at that root succeeds, and every defeater
of \(X\) defeats its superargument \(A\).
Thus it remains to consider \(B\) with a strict root and a nonempty
fallible frontier.

Every construction below uses either the occurrence-sensitive reversed
path, or the two-root contraposition graft, to attack a selected
\(U\in M(B)\).
It has premise and defeasible-rule sets contained in the unions for
\(B,X\), and likewise for last-rule sets.
Under transposition it contains \(X\) as a whole subtree.
Under contraposition it contains whole \(B\), whole \(X\), both, or
neither. A graft containing neither is strict and firm and cannot
be below \(U\); hence it already gives a successful attack.
All other cases are handled by the following finite resource comparisons.

\emph{Last link, defeasible attacker.}
Put \(S=\operatorname{Last}(B)\ne\varnothing\).
If \(X\) is defeasible-rooted with top rule \(t\), the failed-attack
condition gives \(S\triangleleft\{t\}\).
For elitism select \(r\in S\) minimal in \(S\cup\{t\}\).
For democracy every member of \(S\) is below \(t\), so a minimal
member \(r\) of \(S\) has no smaller member in \(S\cup\{t\}\).
Choose a frontier subtree \(U\) with top rule \(r\).
For elitism, the graft's last-rule set is contained in \(S\cup\{t\}\),
and has no element below \(r\). Thus it is not elitistically below
\(\operatorname{Last}(U)=\{r\}\).
For democracy, a transposition graft contains \(t\), which is not
below \(r\). A contraposition graft containing \(B\) contains \(r\);
one containing only \(X\) contains \(t\). Either supplies a member
not below \(r\). In each case \(W\not\prec U\).

If \(X\) is the ordinary premise \([x]\), its last-rule set is empty.
For either set comparison choose \(r\) minimal in \(S\), with
a frontier subtree \(U\) whose top rule is \(r\).
A transposition graft has last-rule set contained in \(S\).
If that set is nonempty, none of its elements is below \(r\);
if it is empty, the empty-left clause applies.
It is therefore not below \(\{r\}\) under either comparison.
For contraposition, a graft containing \(B\) contains \(r\),
and one using only \(X\) has empty last-rule set.
The same noncomparison follows.

\emph{Last link, strict attacker.}
Here \(B\) is plausible and strict.
Its comparison with \(X\) can hold only if \(X=[x]\) is ordinary,
and \(P_B=\operatorname{Prem}_p(B)\triangleleft\{x\}\).
For elitism choose \(u\in P_B\) minimal in \(P_B\cup\{x\}\);
for democracy all members of \(P_B\) are below \(x\), so choose
one minimal in \(P_B\). It is also minimal in their union.
The selected frontier subtree is the ordinary leaf \(U=[u]\).
Every reversed-path graft is strict with ordinary premises in that union.
The minimality witness blocks elitist comparison with \(\{u\}\).
For democracy the retained \(x\) in a transposition graft is not below
\(u\). A contraposition graft retaining \(B\) contains \(u\), and
one retaining only \(X\) contains \(x\); either blocks the comparison.
Empty ordinary support blocks both comparisons.
Thus \(W\not\prec U\) in this case too.

\emph{Weakest link, defeasible-root target \(X\).}
Inspection of the three ordered comparison cases shows that \(B\)
must be defeasible and
\[
 G=\operatorname{Def}(B)\triangleleft
 H=\operatorname{Def}(X),\qquad G,H\ne\varnothing .
 \tag{5}
\]
For elitism choose a minimal \(r\in G\cup H\) that belongs to \(G\),
using the resource lemma. Select an occurrence of \(r\) in \(B\),
and let \(U\) be the frontier subtree above that occurrence.
Then \(U\) has a defeasible root, and
\[
 r\in\operatorname{Def}(U)\subseteq G.
 \]
The graft's defeasible-rule set is contained in \(G\cup H\),
so it is not elitistically below \(\operatorname{Def}(U)\).
Since \(U\) is defeasible, every applicable weakest-link case
requires this rule comparison; its failure gives \(W\not\prec U\).

For democracy choose a maximal \(d\in G\cup H\) lying in \(H\).
Choose any defeasible frontier subtree \(U\) of \(B\), and put
\(\Sigma=\operatorname{Def}(U)\), a nonempty subset of \(G\).
A transposition graft retains \(X\), hence retains \(d\), which
is below no member of \(\Sigma\).
A contraposition graft containing \(B\) has rule set containing
\(\Sigma\), so cannot be democratically below \(\Sigma\).
One containing only \(X\) retains \(d\). In every case the necessary
rule comparison fails, proving \(W\not\prec U\).

\emph{Weakest link, ordinary target \(X=[x]\).}
The failed-attack condition requires \(B\) to be plausible and
\[
 P_B=\operatorname{Prem}_p(B)\triangleleft\{x\}.
 \tag{6}
\]
For elitism choose \(u\in P_B\) minimal in \(P_B\cup\{x\}\).
Select an occurrence of the ordinary leaf \([u]\) in \(B\) and
the frontier subtree \(U\) above it. This \(U\) is either \([u]\)
or a plausible defeasible argument, and
\(u\in\operatorname{Prem}_p(U)\subseteq P_B\).
The graft's ordinary-premise set is contained in \(P_B\cup\{x\}\),
so cannot be elitistically below \(\operatorname{Prem}_p(U)\).
Since \(U\) is plausible, its weakest-link comparison necessarily
uses this premise comparison; hence \(W\not\prec U\).

For democracy, (6) says that every member of \(P_B\) is below \(x\).
Select any ordinary-premise occurrence and its frontier subtree \(U\);
its ordinary-premise set \(\Sigma\) is nonempty and contained in \(P_B\).
A transposition graft retains \(x\), which is below no member of
\(\Sigma\). A contraposition graft containing \(B\) contains
\(\Sigma\), and one containing only \(X\) contains \(x\).
Thus its ordinary-premise set is not democratically below \(\Sigma\),
and again \(W\not\prec U\).

In every case \(W\)'s contradictory attack at \(U\) succeeds.
It defeats \(B\). Every defeater of \(W\) defeats \(B\) or \(X\),
and a defeat of \(X\) propagates to \(A\). This proves (P).
All constructed trees have premises in
\(\operatorname{Prem}(B)\cup\operatorname{Prem}(X)\),
which is contained in \(\operatorname{Prem}(B)\cup\operatorname{Prem}(A)\).
Joint c-consistency therefore admits them in the restricted universe.
None of the selections requires disjoint subtrees: an occurrence can
remain in a sibling, and the resource arguments use set inclusion or
membership rather than a purported deletion.
\end{proof}
\end{annexitem}

\begin{annexitem}{annex:p18.prioritysupport}{thm:p18-priority-supported-union}{The minimal ordinary-premise proof}
\begin{proof}
Let \(S\) be the union of the premise sets of the finitely many defended
arguments. It is finite. Suppose it is c-inconsistent, and choose
an inclusion-minimal c-inconsistent \(S_0\subseteq S\).
The finite set \(O=S_0\cap K_p\) is nonempty, since otherwise
axiom consistency would be violated.
Choose \(u\in O\) minimal under \(<_p\).
By c-classicality there is a strict derivation of a contradictory
of \(u\) from \(S_0\setminus\{u\}\).
It gives a strict argument \(W\) with c-consistent premises, because
that proper subset is c-consistent. Thus \(W\) is admitted.
Its ordinary-premise set is a subset of \(O\setminus\{u\}\),
and no member of that set is below \(u\).
Consequently it is not below \(\{u\}\) for either strict set comparison:
elitism has no witness, democracy has no larger target for any source
member, and empty source support is covered by the empty-left clause.
Both \(W\) and the premise argument \([u]\) are strict, so either link
principle compares exactly those ordinary-premise sets.
Hence \(W\not\prec[u]\), and \(W\) defeats \([u]\).

Each premise argument used in \(W\), and \([u]\), is a subargument
of one of the given defended arguments, so is defended by \(E\).
Every defeat of the strict graft \(W\) lands at one of its ordinary
premise leaves. Thus \(W\) is also defended by \(E\).
But two arguments defended by a defeat-conflict-free set cannot defeat
one another: if \(W D[u]\), defense of \([u]\) supplies \(Z\in E\)
with \(Z D W\), and defense of \(W\) supplies \(Y\in E\) with
\(Y D Z\), contradicting defeat-conflict-freeness.
This contradiction proves the support-union assertion without assuming
any return property.
\end{proof}
\end{annexitem}

\begin{annexitem}{annex:p18.priorityrationality}{thm:p18-priority-rationality}{Full rationality for the corrected priorities}
\begin{proof}
In the unrestricted case, Theorem~\ref{thm:p18-priority-return}
establishes global (R). All remaining hypotheses are those of
Theorem~\ref{thm:p18-rationality-general}, which proves the claimed
semantic equivalence, direct consistency, and closure properties.

Consider now the c-restricted universe. Let \(E\) be defeat-conflict-free
and \(A,B\) be defended by \(E\).
Their joint premise set is c-consistent by
Lemma~\ref{thm:p18-priority-supported-union}.
If \(B T A\), either it succeeds or it fails.
A successful defeat contradicts joint defense, by the two-defender
argument at the end of the support proof.
For a failed attack, the support-qualified return theorem supplies
an admitted \(W\) with \(W D B\) such that every defeater of \(W\)
defeats \(A\) or \(B\).
Defense of \(B\) gives \(Z\in E\) with \(Z D W\).
Thus \(Z D A\) or \(Z D B\). Defense of that argument gives
\(Y\in E\) with \(Y D Z\), again contradicting conflict-freeness.
Therefore jointly defended arguments do not attack.
This proof includes repeated arguments and self-attacks.

It follows that any defeat-admissible set is attack-conflict-free,
and hence attack-admissible. Conversely attack-admissibility implies
defeat-admissibility since every defeat is an attack.
Both definitions of complete extension use the same defense relation,
so their complete families agree. The established abstract semantics
then give equality of grounded, preferred, and stable extensions as well.

Subarguments of an admitted tree have smaller premise sets, so remain
c-consistent and admitted. If the tree is defended, each subargument
is defended by target-local propagation. Thus complete extensions
contain their subarguments.
If finitely many arguments in a complete extension supply the antecedents
of a strict rule, their joint support is c-consistent by the support
lemma. The grafted rule application is admitted, is defended because
every defeater targets an input argument, and is therefore in the
complete extension. Nullary strict arguments have c-consistent empty
support by axiom consistency and are unattacked. This proves strict
closure, including all finite strict derivations.

For direct consistency, take a conflicting pair in an admissible \(E\).
They are jointly defended. If the conflict already gives an attack at
a fallible root, it contradicts the joint-compatibility result above.
Two strict firm arguments cannot have conflicting conclusions by axiom
consistency. Otherwise contrary well-formedness permits the same
orientation used in the proof of
Theorem~\ref{thm:p18-rationality-general}: a fallible strict-root target
has a contradictory opposing conclusion.
Reverse a strict path, or use the contraposition graft, to construct
an admitted argument attacking that target.
The new tree has support in the original pair's c-consistent joint
support. It is defended because every defeater targets one of the
original arguments. It and the target are therefore jointly defended,
contradicting attack compatibility.
This establishes direct consistency for admissible extensions.
For complete extensions, strict closure makes this also indirect
consistency.
Only finite trees and finite support unions were used, so the proof
does not restrict the cardinality of the ambient language or universe.
\end{proof}
\end{annexitem}

\begin{annexitem}{annex:p18.priorityrestricted}{thm:p18-priority-restricted-boundary}{A restricted framework without global return}
\begin{proof}
Let \(L\) be the classical propositional formulas over the distinct atoms
\(a,p,r,\delta\). Conflicts are symmetric classical negation:
\(\alpha,\beta\) contradict when one is the syntactic negation of the
other. Let \(R_s\) consist of every classically valid finite inference,
including zero-antecedent tautologies. Set
\[
 K_n=\{\neg p\},\qquad K_p=\{a,p\},\qquad
 R_d=\{d:a\Rightarrow r\},\qquad n(d)=\delta .
 \tag{7}
\]
Give ordinary premises the total strict priority \(a<_p p\);
the single defeasible rule has its singleton preorder.

Finite strict derivations preserve truth in every Boolean valuation.
Conversely, a classically valid inference from any finite premise set
is itself in \(R_s\).
Thus a finite support is c-consistent exactly when it is classically
satisfiable: an unsatisfiable finite set strictly entails an atom and
its negation, whereas a satisfying valuation excludes such a pair.
Any minimal c-inconsistent set is finite, by finite derivation witnesses.
Removing a member \(\phi\) from it classically entails \(\neg\phi\);
otherwise a valuation satisfying the remainder and \(\phi\) would
satisfy the original set. The corresponding strict rule proves
c-classicality.
Closure under contraposition follows directly: take the finite premises
actually used in a strict derivation of \(\psi\). If \(s\) is among
them, replacing it by \(\neg\psi\) entails \(\neg s\); if it is not,
those premises and \(\neg\psi\) are inconsistent and still entail
\(\neg s\). A finite valid rule supplies the required derivation.
The axiom set \(\{\neg p\}\) is satisfiable, so its strict closure is
consistent. All conflicts are symmetric, so well-formedness holds.

The following admitted arguments exist:
\[
 U=[a],\quad N=[\neg p],\quad V=[U;d],\quad
 B=[V,N;\ r,\neg p\to\neg p],\quad A=[p].
 \tag{8}
\]
The final strict inference is classically valid, although one of its
antecedents is unnecessary for classical consequence. No minimal-support
restriction forbids the displayed tree.
The premise set of \(B\) is \(\{a,\neg p\}\), and that of \(A\)
is \(\{p\}\); each is c-consistent.
The argument \(B\) undermines \(A\).
For last link, \(\{d\}\triangleleft\varnothing\), so \(B\prec A\).
For weakest link both
\(\{a\}\triangleleft\{p\}\) and
\(\{d\}\triangleleft\varnothing\) hold.
These comparisons hold for elitism and democracy alike.
Hence \(B T A\) fails as a defeat under each of the four choices.

We show that no argument defeats \(B\).
Its only fallible targets are its ordinary premise \(a\), the
defeasible conclusion \(r\), and the name \(\delta\) of that rule.
An attacker must therefore conclude respectively
\(\neg a,\neg r,\neg\delta\).
No admitted argument has any of these conclusions.
Indeed, an admitted support is a subset of \(\{a,p,\neg p\}\)
which does not contain both \(p\) and \(\neg p\).
Choose a Boolean valuation with \(a,r,\delta\) true, and choose the
value of \(p\) to satisfy that support.
Every strict rule preserves truth in this valuation; the only defeasible
rule \(a\Rightarrow r\) also preserves truth in this chosen valuation.
Induction on the entire argument tree therefore makes its conclusion
true in that valuation. None of \(\neg a,\neg r,\neg\delta\) is true,
so none can be its conclusion.
Thus \(B\) is unattacked and has no returning defeater \(W D B\).
Global (R) fails.

The combined support \(\{a,p,\neg p\}\) is c-inconsistent, exactly
the case excluded by the support-qualified return.
Moreover \(N\) is an unattacked strict firm defeater of \(A=[p]\).
Consequently \(A\) cannot be defended by any defeat-conflict-free
set: defending it would require defeating the unattacked \(N\).
The failed pair therefore cannot occur as a jointly defended pair
in the rationality proof.
\end{proof}
\end{annexitem}

\begin{annexitem}{annex:p18.priorityworked}{thm:p18-priority-worked}{A full five-argument comparison}
\begin{proof}
Only \(a,b\) are premises. The unique argument for \(u\) uses \(d_1\),
the unique argument for \(t\) then uses \(d_2\), and the unique
argument for \(\neg t\) uses \(d_3\) from \(b\). Neither \(t\) nor
\(\neg t\) feeds another rule. Thus the five displayed arguments are exhaustive.
No argument concludes \(\neg a,\neg b,\neg u\), or a negative rule
name. The only attacks are the opposing defeasible-root rebuttals
\(A T B,B T A\).

Their last-rule sets are \(\{d_2\}\) and \(\{d_3\}\).
Since \(d_3<d_2\), either last-link comparison gives \(B\prec A\).
Only \(A D B\) succeeds.
The arguments \(A_0,B_0,I,A\) are unattacked and must be in every
complete extension; \(B\) cannot be included. The resulting set is
conflict-free and stable, proving uniqueness in all four semantics.

For weakest-link elitism,
\(\{a\}\triangleleft_{\rm E}\{b\}\) and
\(\{d_1,d_2\}\triangleleft_{\rm E}\{d_3\}\), the latter witnessed
by \(d_1<d_3\). Thus \(A\prec B\), while the reverse premise
comparison fails. Only \(B D A\) succeeds, and the same forced-argument
calculation gives the unique extension \(\{A_0,B_0,I,B\}\).

For weakest-link democracy,
\(\{d_1,d_2\}\not\triangleleft_{\rm D}\{d_3\}\), because \(d_2\)
is not below \(d_3\). Thus \(A\not\prec B\).
Although \(\{d_3\}\triangleleft_{\rm D}\{d_1,d_2\}\) is witnessed
by \(d_2\), the premise comparison \(\{b\}\triangleleft_{\rm D}\{a\}\)
fails, so \(B\not\prec A\).
Both rebuttals succeed.
Every complete extension contains \(U_0=\{A_0,B_0,I\}\).
This set attacks neither \(A\) nor \(B\), so defends neither and is
itself complete. Adding either \(A\) or \(B\) makes that argument
defend itself against the other, and yields a stable four-argument
extension. Conflict-freeness prevents adding both.
These are all complete extensions, and their two maximal members
are exactly the preferred and stable extensions; their least member
\(U_0\) is grounded.
In each one-direction case the successful opposite argument is the
return for the failed attack, and every defeat of that returning
argument already defeats the corresponding original target.
\end{proof}
\end{annexitem}
